% Discussion for lab1_item4 -- never overwritten by make_latex.py.

Yes, a directory can be linked using a symbolic link on Linux (\texttt{ln -s})
or a directory junction on Windows (\texttt{mklink /J}); hard links to
directories are forbidden on both operating systems (\texttt{ln} reports
\texttt{hard link not allowed for directory} and \texttt{mklink /H} fails with
\texttt{Access is denied}), and unprivileged \texttt{mklink /D} fails for lack
of symlink privilege.

Prohibiting directory hard links keeps the core directory hierarchy acyclic so
reference-counted garbage collection cannot leak isolated cycles. Allowing
symbolic links or junctions turns the directory namespace into a general graph
in which cycles such as \texttt{parent/child/loop -> ..} can be created.

Recursive traversal tools handle the cycle differently:
\begin{itemize}
  \item On Linux, \texttt{ls -R parent} does not follow symbolic links by
    default, so it lists \texttt{loop} as an entry and terminates without
    recursing. When link-following is enabled, both \texttt{ls -RL}
    (\texttt{not listing already-listed directory}) and \texttt{find -L}
    (\texttt{File system loop detected}) track visited \texttt{(device, inode)}
    pairs and abort the cycle.
  \item On Windows, \texttt{dir /s parent} traverses directory junctions without
    cycle detection, descending repeatedly into
    \texttt{parent\textbackslash child\textbackslash loop\textbackslash child\textbackslash loop\textbackslash\ldots}
    until the Windows path-length limit stops the recursion after
    \winDirLines{} lines of output.
\end{itemize}
